\section{A Unified Taxonomy: The Five Evolution Loops}
\label{sec:taxonomy}

The diversity of self-evolving AI systems---spanning AutoML, neural architecture search, LLM self-improvement, evolutionary computation, and agent frameworks---presents a classification challenge. Each community has developed its own terminology: AutoML researchers speak of ``pipeline synthesis,'' NAS researchers of ``architecture search,'' LLM researchers of ``self-refinement,'' evolutionary computation researchers of ``mutation and selection,'' and agent researchers of ``role adaptation.'' Yet beneath this terminological diversity lies a common structure: iterative cycles of generation, evaluation, memory, and refinement.

In this section, we propose a unified taxonomic framework---the Five Evolution Loops---that captures this common structure. Each loop type represents a fundamental pattern of self-improvement that recurs across domains. Individual systems may employ one or more loop types, and the most powerful systems combine multiple loops into integrated self-evolution pipelines. For each loop, we provide a formal definition, identify representative systems with detailed algorithmic descriptions and experimental data, analyze limitations and failure modes, and examine how loops interact with one another. We conclude with two detailed case studies demonstrating multi-loop composition.

\subsection{Overview of the Five Loops}

We identify five fundamental types of evolutionary loops:

\begin{enumerate}
    \item \textbf{Specification-to-Execution Loop} (\emph{Spec Loop}): Translating high-level goals (typically expressed in natural language) into executable artifacts (ML pipelines, agent workflows, code).
    \item \textbf{Search Loop}: Systematically exploring a space of candidates (architectures, prompts, code, agents, hyperparameters) to find high-performing solutions.
    \item \textbf{Evaluator Loop}: Using automated evaluation mechanisms (tests, benchmarks, verifiers, reward models) to score and select candidate solutions.
    \item \textbf{Reflection Loop}: Analyzing failures and successes, extracting lessons, and incorporating those lessons into future generation or search strategies.
    \item \textbf{Population Loop}: Maintaining a diverse collection of candidates, applying selection, mutation, and recombination operators, and preserving lineage information.
\end{enumerate}

These loops are related but distinct. The Search Loop and Population Loop both explore candidate spaces, but the Search Loop may operate with a single candidate (e.g., hill climbing on prompts) while the Population Loop explicitly maintains multiple simultaneous candidates. The Evaluator Loop and Reflection Loop both process feedback, but the Evaluator Loop produces scalar scores while the Reflection Loop produces structured, language-based analyses. The Specification-to-Execution Loop is unique in its focus on the \emph{translation} from intent to artifact.

Figure~\ref{fig:five_loops} illustrates the relationships between the five loops.

\begin{figure}[t]
\centering
\begin{tikzpicture}[
    node distance=2cm,
    box/.style={draw, rounded corners, minimum width=3cm, minimum height=1cm, align=center, font=\small},
    arrow/.style={->, thick, >=stealth}
]
\node[box, fill=blue!10] (spec) {Specification-to-\\Execution Loop};
\node[box, fill=green!10, right=3cm of spec] (search) {Search Loop};
\node[box, fill=orange!10, below=1.5cm of search] (eval) {Evaluator Loop};
\node[box, fill=red!10, left=3cm of eval] (reflect) {Reflection Loop};
\node[box, fill=purple!10, below=1.5cm of reflect] (pop) {Population Loop};

\draw[arrow] (spec) -- node[above, font=\scriptsize]{specification} (search);
\draw[arrow] (search) -- node[right, font=\scriptsize]{candidate} (eval);
\draw[arrow] (eval) -- node[below, font=\scriptsize]{feedback} (reflect);
\draw[arrow] (reflect) -- node[left, font=\scriptsize]{lessons} (search);
\draw[arrow, dashed] (pop) -- node[left, font=\scriptsize]{diversity} (search);
\draw[arrow, dashed] (pop) -- node[right, font=\scriptsize]{selection} (eval);
\draw[arrow, dashed] (pop) -- node[below, font=\scriptsize]{lineage} (reflect);
\draw[arrow, bend left=20] (eval) to node[right, font=\scriptsize]{scores} (pop);
\end{tikzpicture}
\caption{The Five Evolution Loops and their relationships. The central cycle (Search $\rightarrow$ Evaluate $\rightarrow$ Reflect) is augmented by the Population Loop (maintaining diversity and lineage) and the Specification-to-Execution Loop (translating goals into artifacts).}
\label{fig:five_loops}
\end{figure}


\subsection{Formal Framework}

Before detailing each loop, we establish a formal framework that enables precise comparison. A self-evolving system $\mathcal{S}$ is characterized as a tuple $\mathcal{S} = (\mathcal{G}, \mathcal{E}, \mathcal{M}, \mathcal{R})$ where:

\begin{itemize}
    \item $\mathcal{G}: \mathcal{M} \rightarrow \mathcal{C}$ is a \textbf{generator} mapping memory state to candidate artifacts in candidate space $\mathcal{C}$.
    \item $\mathcal{E}: \mathcal{C} \rightarrow \mathcal{F}$ is an \textbf{evaluator} mapping candidates to feedback signals in feedback space $\mathcal{F}$.
    \item $\mathcal{M}$ is a \textbf{memory} that evolves over time: $\mathcal{M}_{t+1} = \text{Update}(\mathcal{M}_t, c_t, f_t)$.
    \item $\mathcal{R}: (\mathcal{G}, \mathcal{M}) \rightarrow \mathcal{G}'$ is a \textbf{refiner} that updates the generator based on accumulated experience.
\end{itemize}

The five loops operate on different components of this framework:

\begin{table}[t]
\centering
\small
\caption{Mapping of Five Evolution Loops to formal framework components.}
\label{tab:loop_components}
\begin{tabular}{@{}llll@{}}
\toprule
\textbf{Loop} & \textbf{Operates on} & \textbf{Key output} & \textbf{State change} \\
\midrule
Specification & Problem definition & Initial artifact $a_0 \in \mathcal{C}$ & $\emptyset \rightarrow \mathcal{M}_0$ \\
Search & Generator $\mathcal{G}$ & Candidate $c_t$ & $\mathcal{G}_t \rightarrow \mathcal{G}_{t+1}$ \\
Evaluator & Evaluation $f$ & Score $f(c_t)$ & Appends to $\mathcal{M}$ \\
Reflection & Memory $\mathcal{M}$ & Structured feedback $r_t$ & Enriches $\mathcal{M}$ \\
Population & Multiple candidates & Population $P_t$ & $P_t \rightarrow P_{t+1}$ \\
\bottomrule
\end{tabular}
\end{table}


\subsection{Loop 1: Specification-to-Execution}
\label{sec:spec-loop}

\subsubsection{Definition}

The Specification-to-Execution Loop (Spec Loop) transforms a high-level description of a goal---typically expressed in natural language---into an executable artifact. This artifact may be a machine learning pipeline, an agent workflow, a software program, or a configuration file. The Spec Loop is the entry point for self-evolution: it determines what the system will attempt to improve.

Formally, the Spec Loop maps a specification $s \in \mathcal{S}$ to an executable artifact $a \in \mathcal{A}$:
\begin{equation}
    a = \text{SpecLoop}(s) = \text{Plan}(\text{Understand}(s), \text{Decompose}(s), \text{Execute}(\cdot))
\end{equation}

The Spec Loop typically involves three sub-processes: (1)~understanding the specification (parsing the user's intent), (2)~decomposing the specification into sub-tasks (planning), and (3)~executing the sub-tasks (implementation). The key challenge is bridging the semantic gap between natural language, which is ambiguous and context-dependent, and executable artifacts, which must be precise and syntactically correct.

\subsubsection{Representative Systems}

\paragraph{AutoML-GPT.}
AutoML-GPT~\citep{zhang2023automlgpt} is a pioneering system that uses GPT models as the planning layer for automated machine learning. The user provides a task card describing the dataset and objective in natural language, and the system generates a complete ML pipeline including data preprocessing, model selection, training configuration, and evaluation strategy. The system operates through a code-generation approach: it outputs Python code that implements the complete ML pipeline, which is then executed in a sandboxed environment. Results are reported back to the user in natural language. The key limitation is its single-shot nature: it generates a pipeline once without iteratively improving it based on execution results.

\paragraph{AutoML-Agent.}
AutoML-Agent~\citep{chen2024automlagent} extends AutoML-GPT by decomposing the AutoML process into a multi-agent workflow with five specialized agents: (1)~a Prompt Agent that interprets the natural language task specification and extracts structured requirements, (2)~a Manager Agent that coordinates the overall workflow using a 7-state finite state machine, (3)~a Data Agent that retrieves and preprocesses data, (4)~a Model Agent that selects and configures ML models, and (5)~an Operation Agent that handles training, evaluation, and deployment. This decomposition enables more robust and interpretable pipeline generation. AutoML-Agent introduces an Evol-Instruct mechanism for prompt evolution: the system iteratively refines its task understanding by generating variant specifications and testing them against the user's original description. The system employs a 4-way knowledge retrieval mechanism combining web search, local documentation, model cards, and dataset descriptions.

On seven ML task types (classification, regression, clustering, time series, NLP, computer vision, and tabular), AutoML-Agent achieves competitive or superior performance compared to traditional AutoML baselines.

\paragraph{LLM+AutoML Synergy.}
The work on LLM synergies with automated machine learning~\citep{zheng2024llm_automl} demonstrates that LLMs can serve as the bridge between natural language task descriptions and the ML program synthesis process. Rather than treating AutoML as a fixed toolbox, this approach treats it as a program synthesis problem where the LLM proposes and the AutoML system verifies.

\paragraph{AI Scientist.}
The AI Scientist~\citep{lu2024aiscientist} represents the most ambitious Spec Loop system: it translates the specification ``perform novel scientific research'' into an executable research pipeline including idea generation, experimental design, code execution, manuscript writing, and peer review. The AI Scientist-v2 achieved the first AI-generated paper accepted through peer review at a machine learning workshop, at a cost of approximately \$15 per paper. The system scores ideas on novelty, feasibility, and significance: $f(\text{idea}) = \text{novelty} \times \text{feasibility} \times \text{significance}$.

\subsubsection{Limitations and Cross-Loop Interactions}

The Spec Loop faces several fundamental challenges. First, the semantic gap between natural language specifications and executable artifacts is large and domain-dependent: a specification like ``build a good classifier for medical images'' requires substantial domain knowledge to translate into a concrete pipeline with appropriate preprocessing, model architecture, training configuration, and evaluation strategy. The ambiguity of natural language means that different interpretations of the same specification may lead to vastly different artifacts, and there is no automated mechanism for determining which interpretation best matches the user's intent.

Second, most current Spec Loop systems are single-shot: they generate an artifact once without iteratively improving it. The integration of the Spec Loop with the Evaluator and Reflection loops---enabling the system to observe that its generated artifact underperforms and revise the specification or implementation accordingly---remains an open challenge. This integration would require the system to trace performance failures back to specification ambiguities or implementation errors, a form of causal reasoning that current LLMs do not perform reliably.

Third, the verification problem is acute: how can the system determine whether the generated artifact faithfully implements the specification? For formally specified tasks (e.g., mathematical optimization), verification is straightforward. For open-ended tasks (e.g., ``build an effective agent for customer service''), verification requires evaluation against implicit criteria that may not be fully captured in the specification. The AI Scientist addresses this through automated peer review, but the quality of AI-generated reviews remains significantly below human standards.

Fourth, the Spec Loop's performance is fundamentally limited by the LLM's domain knowledge. AutoML-GPT's pipeline generation quality depends on the training data's coverage of ML techniques; the AI Scientist's research quality depends on the LLM's understanding of the relevant scientific literature. This creates a knowledge gap that is difficult to close without explicit knowledge retrieval or domain-specific fine-tuning.

The Spec Loop interacts with other loops in several ways. It feeds into the Search Loop by defining the search space (the artifact type and constraints). It interacts with the Evaluator Loop by providing the success criteria against which candidates are evaluated. It can be enhanced by the Reflection Loop when the system analyzes why its specification-to-execution translation failed and revises its approach. In the most sophisticated systems (AI Scientist), the Spec Loop itself becomes a target of evolution: the system learns to write better specifications over time, improving both its understanding of user intent and its ability to generate high-quality artifacts from specifications.


\subsection{Loop 2: Search}
\label{sec:search-loop}

\subsubsection{Definition}

The Search Loop systematically explores a space of candidate solutions to find high-performing configurations. Formally, the Search Loop iteratively generates candidates:
\begin{align}
    c_{t+1} &= \text{SearchStep}(c_t, \mathcal{H}_t, \mathcal{F}) \\
    \mathcal{H}_{t+1} &= \mathcal{H}_t \cup \{(c_t, f(c_t))\}
\end{align}
where $c_t$ is the current candidate, $\mathcal{H}_t$ is the history of past candidates and scores, $\mathcal{F}$ is the search strategy, and $f$ is the evaluation function.

\subsubsection{Search Strategies}

\paragraph{LLM-based optimization (OPRO).}
OPRO~\citep{yang2023opro} uses an LLM as a black-box optimizer. The optimization problem is described in natural language, and the LLM reads a history of past solutions paired with their scores to propose new candidates. The meta-prompt structure encodes the optimization trajectory:
\begin{verbatim}
Instruction Score: 45
"Let's think step by step"
Instruction Score: 62
"Think about this problem carefully"
Instruction Score: 78
"Break down the problem into parts"
Instruction Score: [to be filled]
\end{verbatim}

On GSM8K, OPRO discovers prompts achieving 80\% accuracy with PaLM 2-L (vs.\ 57\% default), a 23pp improvement. On BIG-Bench Hard, improvements reach 25pp. The method also applies to linear regression and traveling salesman problems, demonstrating universality.

\paragraph{Evolutionary search (AlphaEvolve, LLaMEA, ReEvo).}
AlphaEvolve~\citep{novikov2025alphaevolve} uses a dual-model strategy: Gemini Flash for breadth exploration and Gemini Pro for depth exploration. Key results include 48 multiplications for $4\times4$ complex matrices (first Strassen improvement in 56 years), 0.7\% worldwide compute recovery in data center scheduling, 23\% FlashAttention speedup, and 32\% LLM training attention speedup.

LLaMEA~\citep{van2024llaMEA} evolves optimization algorithms at the meta-level, using SEARCH/REPLACE diff mode and SMAC in-loop hyperparameter optimization. It won the GECCO 2025 Silver Humies award.

ReEvo~\citep{ye2024reevo} (NeurIPS 2024) introduces ``Language Hyper-Heuristics'' combining short-term and long-term reflection. On TSP with 20 nodes, ReEvo achieves near-optimal solutions (gap $<$ 0.5\%) while discovering novel heuristics.

\paragraph{Architecture search (NAS with LLMs).}
EvoPrompting~\citep{chen2024evoprompting} encodes GNN architectures as Python code and evolves them via LLM. On the CLRS benchmark, EvoPrompting discovers architectures outperforming human-designed GNNs. NADER~\citep{nasr2024nader} decomposes architecture search into multi-agent collaboration. EvoPrompt~\citep{guo2023evoprompt} treats prompts as evolvable organisms, achieving up to 25\% improvement on BBH and 15\% on GSM8K across 31 datasets.

\paragraph{Agent architecture search (ADAS, DGM).}
ADAS~\citep{hu2024adas} searches over the space of all valid Python agent implementations (Turing-complete search space) using a three-layer generation-reflection-debug process with dynamic \texttt{exec()} evaluation and bootstrap confidence intervals. Starting from five seed architectures, ADAS discovered novel architectures that outperform hand-designed baselines on ARC and GFootball. Critically, architectures discovered by GPT-3.5 transfer to GPT-4, demonstrating generalizable architectural principles.

The Darwin G\"odel Machine~\citep{zhang2025dgm} maintains an open-ended archive of agent implementations with sampling proportional to $f(a_i) \times \text{diversity\_bonus}(a_i, A)$. On SWE-bench: 20.0\% $\rightarrow$ 50.0\%; on Polyglot: 14.2\% $\rightarrow$ 30.7\%.

\subsubsection{Analysis}

The key trend is LLMs as universal search operators. Traditional algorithms require hand-coded operators for each space; LLM-based operators apply to any artifact expressible in natural language or code. A fundamental limitation is the \emph{search space specification problem}: defining the space to be both expressive enough and tractable. ADAS addresses this with Turing-complete Python code; OPRO uses natural language for flexibility.

\subsubsection{Scaling Properties}

Several important scaling relationships emerge from the empirical results across search-based systems:

\begin{itemize}
    \item \textbf{Search budget vs.\ quality}: AlphaEvolve's discoveries required evaluating thousands of candidate programs. The 48-multiplication Strassen improvement emerged after maintaining a MAP-Elites archive across multiple problem sizes and running parallel evaluations for extended periods. This suggests that for hard discovery problems, search budget is a primary determinant of success.

    \item \textbf{History length vs.\ optimization quality}: OPRO's performance improves with the length of the optimization history provided to the LLM, but with diminishing returns. Providing more than 20 past solutions yields marginal improvements, suggesting that the LLM's ability to extract patterns from optimization trajectories saturates at moderate history lengths.

    \item \textbf{Model capability vs.\ search effectiveness}: ADAS experiments show that the quality of discovered architectures correlates with the search-time model's capability: GPT-4 discovers better architectures than GPT-3.5, but the gap is smaller than the capability gap between the models, suggesting that the search framework provides a significant boost that partially compensates for model limitations.

    \item \textbf{Transfer across domains}: Architectures discovered by ADAS on one task transfer to related tasks, and architectures designed by one model transfer to execution by a different model. This transfer property is crucial for practical deployment: it means that expensive architecture search can be amortized across multiple applications.
\end{itemize}

\subsubsection{Failure Modes}

Search-based systems exhibit several characteristic failure modes:

\begin{enumerate}
    \item \textbf{Mode collapse}: The search converges to a narrow region of the solution space, producing variants of the same approach. Population-based methods (MAP-Elites, island models) address this explicitly through diversity maintenance, but single-trajectory methods (OPRO, Self-Refine) are susceptible.

    \item \textbf{Evaluator exploitation}: When the evaluation function has exploitable weaknesses, the search process will discover candidates that achieve high scores by gaming the evaluator rather than solving the underlying problem. This is particularly acute with learned reward models.

    \item \textbf{Search space mismatch}: When the search space is too narrow (excluding good solutions) or too broad (containing too many poor solutions), search efficiency degrades. The choice of representation (code, natural language, structured configurations) determines the effective search space.

    \item \textbf{Computational cost}: For Turing-complete search spaces (ADAS), the number of possible candidates is unbounded, and the cost of evaluating each candidate (running the agent on benchmark tasks) is substantial. Budget management---how to allocate a fixed compute budget across breadth exploration, depth refinement, and evaluation---is a key practical challenge.
\end{enumerate}


\subsection{Loop 3: Evaluator}
\label{sec:evaluator-loop}

\subsubsection{Definition}

The Evaluator Loop provides selection pressure. Formally, it applies an evaluation function:
\begin{equation}
    f: \mathcal{C} \rightarrow \mathbb{R}^k
\end{equation}
where $\mathcal{C}$ is the candidate space and $\mathbb{R}^k$ is a vector of potentially multi-objective scores.

We identify an ``evaluation ladder'' of increasing reliability: human evaluation (gold standard, unscalable) $\rightarrow$ reward models (scalable, biased) $\rightarrow$ LLM-based evaluation (flexible, limited) $\rightarrow$ benchmark suites (reproducible, narrow) $\rightarrow$ execution-based (deterministic, domain-limited) $\rightarrow$ mathematical verification (strongest guarantee, narrowest applicability).

\subsubsection{Evaluation Mechanisms}

\paragraph{Execution feedback.}
Self-Debug~\citep{chen2024selfdebug} uses execution traces as evaluation signals. SelfEvolve~\citep{jiang2024selfevolve} extends this with self-generated API documentation, achieving 85.98\% HumanEval pass@1 (from 74.39\% baseline). The mechanism is deterministic and reliable but requires well-defined test suites.

\paragraph{Mathematical verification.}
FunSearch~\citep{romera2024funsearch} uses 10 islands with 15 samplers and 140 evaluators running in parallel, with signature-level clustering for diversity. AlphaEvolve~\citep{novikov2025alphaevolve} uses multi-dimensional scoring: correctness, efficiency, and novelty. Mathematical verification provides the strongest guarantee---provable correctness---but is applicable only to formally specified domains.

\paragraph{Reward models and AI feedback.}
ReST~\citep{gulcehre2023rest} uses learned reward models to filter generated samples. Constitutional AI~\citep{bai2022constitutional} uses principle-based AI evaluation. Reward models are flexible (can evaluate open-ended outputs) but less reliable (susceptible to reward hacking).

\paragraph{Benchmark suites.}
SWE-bench~\citep{jimenez2024swebench} evaluates on real GitHub issues from 12 Python repositories. SWE-Agent~\citep{yang2024sweagent} (ICLR 2025 Oral) introduces the Agent-Computer Interface. OSWorld~\citep{xie2024osworld} tests OS-level tasks; WebArena~\citep{zhou2024webarena} tests web interaction. These provide realistic multi-step evaluation but may not generalize to deployment.

\paragraph{Quality-diversity metrics.}
MAP-Elites~\citep{mouret2015mapelites} evaluates on both quality and diversity. The archive $\mathcal{A}: \mathcal{B} \rightarrow \mathcal{C}$ stores the best candidate for each behavioral niche. A new candidate $c$ is added iff $f(c) > f(\mathcal{A}(b(c)))$. This prevents premature convergence and maps the full solution landscape.

\subsubsection{Analysis}

The Evaluator Loop is the bottleneck of self-evolving systems. Three manifestations: (1)~many real-world tasks lack automated evaluators; (2)~evaluators may be gameable; (3)~multi-objective evaluation requires trade-off specification. The most promising direction combines multiple mechanisms: execution for verifiable properties, reward models for subjective quality, quality-diversity for exploration. ReVeal~\citep{jin2026reveal} co-evolves the code generator and test generator, creating adversarial evaluation dynamics.

\subsubsection{The Evaluation Quality-Evolution Quality Connection}

The relationship between evaluation quality and evolution quality is not linear but follows a threshold pattern. Below a critical evaluation quality threshold, evolution fails entirely: the system optimizes for evaluator quirks rather than genuine improvement. Above the threshold, evolution quality improves approximately linearly with evaluation quality, with diminishing returns at very high evaluation quality. This threshold effect explains why systems with perfect evaluators (FunSearch, AlphaEvolve) can make discoveries that escape systems with noisy evaluators (RLHF-trained models): the perfect evaluator ensures that every improvement is genuine, not an artifact of evaluation noise.

The threshold effect has practical implications for system design. It suggests that investing in evaluation quality---better test suites, more comprehensive benchmarks, formal verification---has higher returns than investing in search or generation capability. A system with a mediocre generator and a perfect evaluator will outperform a system with a perfect generator and a mediocre evaluator, because the former can evolve to improve its generation while the latter cannot distinguish genuine improvements from evaluator exploitation.

\subsubsection{Emerging Evaluation Paradigms}

Several emerging paradigms promise to address the evaluation bottleneck:

\paragraph{Co-evolved evaluation (ReVeal).}
ReVeal~\citep{jin2026reveal} co-evolves code generation and test generation, using multi-turn RL to optimize both simultaneously. The TAPO (Turn-level Adaptive Policy Optimization) framework assigns per-turn credit to both code and test generation turns. This adversarial co-evolution creates a ``arms race'' where both components continuously improve: better tests drive better code generation, and better code generation exposes weaknesses in existing tests.

\paragraph{Self-play evaluation (Absolute Zero).}
Absolute Zero~\citep{zhao2025absolutezero} uses self-play for evaluation: the task proposer and task solver are the same model, and the task difficulty automatically adapts to the solver's capability. This eliminates the need for external evaluation datasets entirely, replacing them with self-generated evaluation that is always calibrated to the system's current capability level.

\paragraph{Execution-guided evaluation at scale.}
Systems like OpenEvolve~\citep{openevolve2024} and OpenTreeSearch~\citep{opentreesearch2025} demonstrate that execution-based evaluation can scale to complex real-world problems. OpenTreeSearch uses PUCT (Predictor + Upper Confidence Bound for Trees) tree search with a single hyperparameter replacing 9+ island model parameters, achieving Circle Packing $n=26$: 2.632 (vs.\ AlphaEvolve's 2.635) with dramatically simpler configuration.


\subsection{Loop 4: Reflection}
\label{sec:reflection-loop}

\subsubsection{Definition}

The Reflection Loop converts evaluation outcomes into structured, actionable knowledge:
\begin{equation}
    r_t = \text{Reflect}(c_t, f(c_t), \mathcal{R}_{<t})
\end{equation}
The reflection $r_t$ is stored in memory and guides future generation. While evaluation produces scalars, reflection produces structured natural language analyses identifying specific failure modes and proposing targeted fixes.

\subsubsection{Reflection Mechanisms}

\paragraph{Verbal reinforcement learning (Reflexion).}
Reflexion~\citep{shinn2023reflexion} uses three models: Actor ($M_a$), Evaluator ($M_e$), and Self-Reflection ($M_r$). The loop for each episode $e$:
\begin{align}
    a_e &= M_a(\text{task}, M_{1:e-1}) \\
    r_e &= M_e(a_e, \text{task}) \\
    f_e &= M_r(a_e, r_e, \text{task}) \quad \text{if failure} \\
    M_e &= M_{e-1} \cup \{f_e\}
\end{align}
The key insight: natural language provides a richer signal than scalar rewards---telling the system \emph{why} it failed, not just \emph{that} it failed. Reflexion achieved 91\% HumanEval pass@1 with GPT-3.5 (vs.\ GPT-4's 80\%) and 97\% AlfWorld success (vs.\ 75\% ReAct).

\paragraph{Iterative self-refinement (Self-Refine).}
Self-Refine~\citep{madaan2023selfrefine}: Generate $y_0$ $\rightarrow$ Critique $\text{fb}_t$ $\rightarrow$ Refine $y_{t+1}$. Same model for all three roles. Achieves ~20\% improvement across seven tasks: code optimization +28\%, math reasoning 56\%$\rightarrow$67\%, constrained generation +13pp. No training or external feedback required.

\paragraph{Language backpropagation (Agent Symbolic Learning).}
Agent Symbolic Learning~\citep{zhou2024agentsymbolic} treats agents as symbolic networks with learnable ``weights'' (prompts, tools, pipelines). Four stages:
\begin{enumerate}
    \item Forward pass: Execute agent, producing trajectory $\tau$.
    \item Language loss: $\mathcal{L}_{\text{lang}} = \text{LLM}(\mathcal{P}_{\text{loss}}(\tau))$.
    \item Language gradient: $\nabla_{\text{lang}}^n = \text{LLM}(\mathcal{P}_{\text{gradient}}(\nabla_{\text{lang}}^{n+1}, \text{info}_n, \mathcal{L}_{\text{lang}}))$.
    \item Weight update: Apply language optimizers to prompts, tools, and pipelines.
\end{enumerate}
Results: HotPotQA F1=39.1 (vs.\ 36.2 baseline), MATH 60.7\% (vs.\ 56.0\%), HumanEval 73.2\% (vs.\ 68.3\%).

\paragraph{Self-play reasoning (Absolute Zero).}
Absolute Zero~\citep{zhao2025absolutezero} proposes tasks and solves them using zero external data. Curriculum reward: $R_{\text{proposer}} \propto \text{difficulty}(t) \times \text{solvability}(t)$. AZR-Coder-7B achieves SOTA at 7B level on HumanEval+, MBPP+, LiveCodeBench.

\subsubsection{Analysis}

The Reflection Loop is uniquely LLM-native---no classical analogue exists. The key tension is specificity vs.\ generality: specific reflections are actionable but may not generalize; general reflections transfer but require additional reasoning. The most effective systems combine both. A fundamental limitation is the lack of convergence guarantees: unlike gradient-based optimization, reflection-based improvement has no theoretical bounds on convergence rate or quality.

\subsubsection{Reflection Quality Taxonomy}

We propose a four-level taxonomy of reflection quality, based on the depth and transferability of the analysis:

\begin{enumerate}
    \item \textbf{Level 0 --- Binary feedback}: ``This solution is incorrect.'' Provides no directional information. Corresponds to pure Evaluator Loop output without reflection.

    \item \textbf{Level 1 --- Surface diagnosis}: ``The error is on line 42: variable \texttt{x} should be initialized to zero.'' Identifies the specific failure point but provides no general insight. Self-Debug operates primarily at this level through execution trace analysis.

    \item \textbf{Level 2 --- Pattern identification}: ``This system tends to forget edge cases involving empty inputs.'' Identifies recurring patterns across failures, enabling proactive avoidance. Reflexion achieves this through accumulated verbal memory.

    \item \textbf{Level 3 --- Principle extraction}: ``Divide-and-conquer strategies work better than brute-force approaches for this class of problems.'' Extracts transferable principles that apply across tasks. Agent Symbolic Learning's language backpropagation aims for this level by propagating gradients through the full agent computational graph.
\end{enumerate}

Higher reflection levels require more sophisticated reasoning but provide more valuable feedback. Current systems operate primarily at Levels 0--2; achieving reliable Level 3 reflection remains an open challenge.

\subsubsection{Memory Management for Reflection}

A practical challenge for reflection-based systems is memory management. Reflexion's linguistic memory buffer grows unboundedly: each failed episode adds a new reflection entry, and the buffer eventually exceeds the model's context window. Current approaches to this problem include:

\begin{itemize}
    \item \textbf{Fixed-size buffers with summarization}: Periodically summarize older reflections into compact representations, preserving key insights while reducing token count. However, summarization may lose important details.

    \item \textbf{Relevance-based retrieval}: Instead of including all past reflections in the context, retrieve only those relevant to the current task. This requires a retrieval mechanism that can identify which past experiences are applicable to the current situation.

    \item \textbf{Hierarchical memory}: Store reflections at multiple levels of abstraction---specific episode-level reflections, intermediate pattern-level summaries, and high-level principle extractions. Navigate the hierarchy based on the current need.

    \item \textbf{Parametric integration}: Move frequently-used reflections into the model's parameters through fine-tuning, freeing context window space for new reflections. STaR and RISE implicitly perform this integration by fine-tuning on self-generated reasoning chains.
\end{itemize}

The choice of memory management strategy affects both the efficiency and the capability of the reflection system. Systems that accumulate rich histories (Reflexion) can provide more context-specific guidance but face scalability challenges. Systems that integrate reflections into parameters (STaR) achieve persistent improvement but may lose the flexibility of context-based adaptation.


\subsection{Loop 5: Population}
\label{sec:population-loop}

\subsubsection{Definition}

The Population Loop maintains diverse candidates and applies evolutionary operators:
\begin{align}
    P' &= \text{Select}(P_t, f) \\
    P'' &= \text{Vary}(P', \text{mutate}, \text{recombine}) \\
    P_{t+1} &= \text{Replace}(P_t, P'', f)
\end{align}
where variation operators are typically LLM-implemented.

\subsubsection{Population Strategies}

\paragraph{MAP-Elites quality-diversity search.}
MAP-Elites~\citep{mouret2015mapelites} maintains a map $\mathcal{A}: \mathcal{B} \rightarrow \mathcal{C}$ across behavioral dimensions. AlphaEvolve's application produced the 48-multiplication algorithm for $4\times4$ complex matrices, using stepping-stone algorithms as foundations for breakthroughs.

\paragraph{Island models (FunSearch).}
FunSearch~\citep{romera2024funsearch} uses 10 islands, 15 samplers, 140 evaluators in parallel. Migration between islands shares strategies while geographic isolation maintains diversity. Signature-level clustering groups programs by behavioral patterns.

\paragraph{Evolving optimizers (LLaMEA).}
LLaMEA~\citep{van2024llaMEA} evolves optimization algorithms at the meta-level with nested SMAC hyperparameter optimization---``improving the improvers.'' Won GECCO 2025 Silver Humies.

\paragraph{Open-ended archives (Darwin G\"odel Machine).}
DGM~\citep{zhang2025dgm} maintains a growing archive with sampling $\propto f(a_i) \times \text{diversity\_bonus}(a_i, A)$. Both parent and improved child are retained---parents as stepping stones. SWE-bench: 20.0\%$\rightarrow$50.0\%; Polyglot: 14.2\%$\rightarrow$30.7\%.

\paragraph{Multi-agent population dynamics.}
AgentEvolver~\citep{jiang2024agentevolver} evolves agent tool usage. EvoAgent~\citep{yuan2024evoagent} evolves agent configurations. EvoAgentX~\citep{evoagentx2024} implements four optimizers (TextGrad, MIPRO, AFlow, EvoPrompt) for agent self-evolution.

\subsubsection{Analysis}

The LLM era transformed the Population Loop: universal variation operators and informed variation from candidate histories. Key challenges: computational cost (mitigated by parallel evaluation) and behavioral descriptor specification for quality-diversity methods. Open-ended archives (DGM) represent a promising direction, preserving full evolutionary history as stepping stones for future discovery.

\subsubsection{Comparative Analysis of Population Strategies}

Table~\ref{tab:population_comparison} compares the population strategies discussed above across several dimensions.

\begin{table}[t]
\centering
\small
\caption{Comparison of population strategies for self-evolving systems.}
\label{tab:population_comparison}
\begin{tabular}{@{}lccccc@{}}
\toprule
\textbf{Strategy} & \textbf{Diversity} & \textbf{Scalability} & \textbf{Stepping Stones} & \textbf{Representative} \\
\midrule
MAP-Elites & High & Medium & Implicit & AlphaEvolve \\
Island Model & High & High & Implicit & FunSearch \\
Meta-Evolution & Medium & Medium & Explicit & LLaMEA \\
Open-ended Archive & Very High & Low-Medium & Explicit & DGM \\
Multi-agent Pop. & Medium & Medium & No & EvoAgent \\
\bottomrule
\end{tabular}
\end{table}

The choice of population strategy depends on the target domain and available compute. For mathematical discovery (AlphaEvolve, FunSearch), where candidate evaluation is expensive but parallelizable, island models with quality-diversity selection are most effective. For agent architecture evolution (DGM), where the search space is vast and stepping stones are critical, open-ended archives provide the best exploration-exploitation balance. For meta-optimization (LLaMEA), where the population consists of algorithms rather than solutions, smaller populations with deeper evaluation per individual are preferred.

A critical insight from DGM's open-ended archive is the importance of \emph{stepping stones}: intermediate solutions that are not optimal for the target task but provide the foundation for discovering superior solutions. On SWE-bench, DGM's archive contained agents that solved only 25\% of tasks but served as essential stepping stones for agents that eventually solved 50\%. This suggests that population management should preserve sub-optimal candidates when they represent qualitatively different approaches, not just when they achieve high scores.


\subsection{Cross-Domain Mapping}
\label{sec:cross-domain}

Table~\ref{tab:cross_domain} maps the five loops across five research domains.

\begin{table}[t]
\centering
\small
\caption{Cross-domain mapping of the Five Evolution Loops. Each cell lists representative systems and the loop's role.}
\label{tab:cross_domain}
\resizebox{\textwidth}{!}{%
\begin{tabular}{lccccc}
\toprule
& \textbf{Spec Loop} & \textbf{Search Loop} & \textbf{Evaluator Loop} & \textbf{Reflection Loop} & \textbf{Population Loop} \\
\midrule
\textbf{AutoML + LLM} &
AutoML-GPT, AutoML-Agent &
Pipeline config search &
Dataset metrics, deployment &
Log analysis &
Ensemble methods \\[4pt]

\textbf{NAS + LLM} &
Architecture specification &
EvoPrompting, LLMatic, NADER &
NAS-Bench accuracy &
Design rationale &
MAP-Elites (LLMatic) \\[4pt]

\textbf{LLM Self-Improve} &
Task description &
Prompt optimization &
Task performance &
Self-Refine, Reflexion, STaR &
SPIN self-play \\[4pt]

\textbf{Evo Comp + LLM} &
Objective function &
OPRO, ReEvo, LLaMEA &
Fitness function &
ReEvo reflective search &
AlphaEvolve, FunSearch \\[4pt]

\textbf{Agent Frameworks} &
Goal/task input &
ADAS, DGM &
SWE-bench, task completion &
Agent Symbolic Learning &
EvoAgent, AgentEvolver \\
\bottomrule
\end{tabular}
}
\end{table}

Key patterns:
\begin{enumerate}
    \item \textbf{Reflection is uniquely LLM-native}. No classical analogue exists. Verbal RL and language backpropagation are fundamentally enabled by LLM language capabilities.
    \item \textbf{Population is underrepresented in agents}. Most agent frameworks use single/fixed agents. Population-level agent search (EvoAgent, DGM) is recent but promising given the vast architecture space.
    \item \textbf{Evaluation is the universal bottleneck}. Across all domains, evolution quality is bounded by evaluation quality---whether dataset metrics, benchmark scores, or reward model reliability.
    \item \textbf{Loop composition varies by maturity}. AutoML: Spec+Evaluator. NAS: +Search. Evolutionary: +Population. LLM self-improvement: +Reflection. Advanced systems combine all five.
\end{enumerate}


\subsection{Loop Composition Patterns}
\label{sec:composition}

We identify three composition patterns:

\paragraph{Linear pipeline.}
Sequential: Specification $\rightarrow$ Search $\rightarrow$ Evaluate $\rightarrow$ Reflect $\rightarrow$ Repeat. Simple but may miss cross-loop optimization. Exemplified by AutoML-Agent.

\paragraph{Nested loops.}
One loop inside another. AlphaEvolve nests Search (LLM mutations) inside Population (MAP-Elites) with Evaluator (automated scoring) and Reflection (evolution context). DGM adds meta-loop depth. Trades capability for computational cost.

\paragraph{Parallel loops.}
Multiple loops running concurrently. Agent Symbolic Learning runs forward execution, loss computation, and backpropagation in parallel. EvoAgentX runs four optimizers simultaneously. Maximizes throughput but requires coordination.

\paragraph{Analysis.}
The choice of composition pattern has significant design implications. Linear pipelines are simple but limited. Nested loops enable hierarchical search but increase cost. Parallel loops maximize throughput but risk conflicting updates. The most powerful systems (AlphaEvolve, DGM) use nested composition with multiple loops interacting at different levels of abstraction.


\subsection{Case Studies in Multi-Loop Composition}
\label{sec:case-studies}

\subsubsection{Case Study 1: AlphaEvolve}

AlphaEvolve~\citep{novikov2025alphaevolve} is the most comprehensive multi-loop system, employing four loops:

\begin{itemize}
    \item \textbf{Search}: Gemini Flash (breadth) and Gemini Pro (depth) propose program mutations.
    \item \textbf{Evaluator}: Automated scoring for correctness, efficiency, and novelty.
    \item \textbf{Population}: MAP-Elites maintains diverse candidates across behavioral dimensions.
    \item \textbf{Reflection}: Evolution context provides implicit learning from mutation history.
\end{itemize}

The missing loop is Specification-to-Execution: AlphaEvolve requires human-specified objectives and does not translate natural language goals into evaluation functions. This is an important gap for future work.

\paragraph{Key results.}
$4\times4$ complex matrix multiplication: 48 multiplications (first Strassen improvement since 1969). Data center scheduling: 0.7\% worldwide compute recovered. FlashAttention: 23\% speedup. LLM training attention: 32\% speedup. New mathematical results in combinatorics and number theory.

\paragraph{Analysis.}
AlphaEvolve demonstrates that four-loop composition produces capabilities no single loop achieves. The Population Loop provides diversity; Search provides exploration; Evaluator provides selection pressure; Reflection provides trajectory learning. The success in mathematical discovery---requiring both creativity and rigor---validates multi-loop composition for challenging intellectual tasks.

\subsubsection{Case Study 2: Reflexion}

Reflexion~\citep{shinn2023reflexion} provides a focused three-loop composition:

\begin{itemize}
    \item \textbf{Search}: Actor generates candidate actions guided by accumulated memory.
    \item \textbf{Evaluator}: Binary/scalar evaluator triggers reflection on failure.
    \item \textbf{Reflection}: Self-reflection generates verbal feedback stored in persistent memory.
\end{itemize}

Missing: Population (single trajectory) and Specification (fixed task format).

\paragraph{Key results.}
HumanEval: 91\% with GPT-3.5 (vs.\ 80\% GPT-4). AlfWorld: 97\% (vs.\ 75\% ReAct). Performance scales monotonically with reflection episodes, diminishing after ~10 episodes.

\paragraph{Comparison with AlphaEvolve.}
Where AlphaEvolve uses four loops with substantial infrastructure (dual Gemini models, MAP-Elites, distributed evaluation), Reflexion achieves strong results with three loops and a single LLM. The comparison illustrates a fundamental tradeoff: Reflexion trades generality for efficiency; AlphaEvolve trades efficiency for generality. The choice between focused and comprehensive loop composition depends on the target domain, available compute, and required capability breadth.


\subsection{Input/Output/State Tables}

For reference, Table~\ref{tab:io_tables} summarizes the input, output, state, and key algorithms for each loop.

\begin{table}[t]
\centering
\small
\caption{Input/Output/State summary for each evolution loop.}
\label{tab:io_tables}
\resizebox{\textwidth}{!}{%
\begin{tabular}{@{}lllll@{}}
\toprule
\textbf{Loop} & \textbf{Input} & \textbf{Output} & \textbf{State} & \textbf{Key Algorithms} \\
\midrule
Specification & Natural language goal & Executable artifact & Task decomposition & Planning, code generation \\
Search & Candidate space, history & New candidate $c_{t+1}$ & History $\mathcal{H}_t$ & OPRO, ADAS, EvoPrompt \\
Evaluator & Candidate $c$ & Score $f(c) \in \mathbb{R}^k$ & Evaluation records & Tests, verification, MAP-Elites \\
Reflection & Candidate, score, past reflections & Structured feedback $r_t$ & Reflection memory $\mathcal{R}$ & Reflexion, Self-Refine, TextGrad \\
Population & Multiple candidates & Evolved population $P_{t+1}$ & Population archive & GA, MAP-Elites, island models \\
\bottomrule
\end{tabular}
}
\end{table}


\subsection{Summary}

The Five Evolution Loops taxonomy provides a unified framework for understanding AI self-evolution. By identifying five fundamental evolutionary cycles and mapping their occurrence across five research domains, we establish a common language enabling systematic comparison and cross-pollination. The key insight: systems from different communities are converging on the same fundamental mechanisms, differing primarily in the artifacts they evolve and the search spaces they explore.

Our case studies demonstrate that the most powerful systems combine multiple loops. AlphaEvolve's four-loop architecture achieves algorithm discovery; Reflexion's three-loop architecture achieves strong results efficiently. The composition pattern is itself a design decision affecting capability, cost, and debuggability.

In the following sections, we apply this taxonomy to analyze specific systems and methods in detail, beginning with self-improvement methods (Section~\ref{sec:selfimprovement}), which most heavily employ the Reflection Loop and Evaluator Loop.
